\SourceLeaf{243}%
dum, cuius dimensio una pro evanescente habetur, flexile quidem, sed non extensibile, qualitates superficiei partim a forma pendent, in quam illa reducta concipitur, partim absolutae sunt, atque invariatae manent, in quamcunque formam illa flectatur. Ad has posteriores, quarum investigatio campum geometriae novum fertilemque aperit, referendae sunt mensura curvaturae atque curvatura integra eo sensu, quo hae expressiones a nobis accipiuntur; porro huc pertinet doctrina de lineis brevissimis, pluraque alia, de quibus in posterum agere nobis reservamus. In hoc considerationis modo superficies plana atque superficies in planum explicabilis, e. g. cylindrica, conica etc. tamquam essentialiter identicae spectantur, modusque genuinus indolem superficiei ita consideratae generaliter exprimendi semper innititur formulae \(\BareRoot{(E\mathrm{d}p^2+2F\mathrm{d}p\mathord{.}\mathrm{d}q+G\mathrm{d}q^2)}\), quae nexum elementi cum duabus indeterminatis \(p,q\) sistit. Sed antequam hoc argumentum ulterius prosequamur, principia theoriae linearum brevissimarum in superficie curva data praemittere oportet.
\Article{14}
Indoles lineae curvae in spatio generaliter ita datur, ut coordinatae \(x,y,z\) singulis illius punctis respondentes exhibeantur in forma functionum unius variabilis, quam per \(w\) denotabimus. Longitudo talis lineae a puncto initiali arbitrario usque ad punctum, cuius coordinatae sunt \(x,y,z\), exprimitur per integrale
\[
\int\mathrm{d}w\mathord{.}\BareRoot{((\tfrac{\mathrm{d}x}{\mathrm{d}w})^2+(\tfrac{\mathrm{d}y}{\mathrm{d}w})^2+(\tfrac{\mathrm{d}z}{\mathrm{d}w})^2)}
\]
Si supponimus, situm lineae curvae variationem infinite parvam pati, ita ut coordinatae singulorum punctorum accipiant variationes \(\delta x,\delta y,\delta z\), variatio totius longitudinis invenitur
\[
=\int\frac{\mathrm{d}x\mathord{.}\mathrm{d}\delta x+\mathrm{d}y\mathord{.}\mathrm{d}\delta y+\mathrm{d}z\mathord{.}\mathrm{d}\delta z}{\BareRoot{(\mathrm{d}x^2+\mathrm{d}y^2+\mathrm{d}z^2)}}
\]
quam expressionem in hanc formam transmutamus:
\[
\begin{aligned}
&\frac{\mathrm{d}x\mathord{.}\delta x+\mathrm{d}y\mathord{.}\delta y+\mathrm{d}z\mathord{.}\delta z}{\BareRoot{(\mathrm{d}x^2+\mathrm{d}y^2+\mathrm{d}z^2)}}\\
&-\int(\delta x\mathord{.}\mathrm{d}\frac{\mathrm{d}x}{\BareRoot{(\mathrm{d}x^2+\mathrm{d}y^2+\mathrm{d}z^2)}}+\delta y\mathord{.}\mathrm{d}\frac{\mathrm{d}y}{\BareRoot{(\mathrm{d}x^2+\mathrm{d}y^2+\mathrm{d}z^2)}}\\
&\quad+\delta z\mathord{.}\mathrm{d}\frac{\mathrm{d}z}{\BareRoot{(\mathrm{d}x^2+\mathrm{d}y^2+\mathrm{d}z^2)}})
\end{aligned}
\]
In casu eo, ubi linea est brevissima inter puncta sua extrema, constat. ea, quae hic sub signo integrali sunt, evanescere debere. Quatenus linea esse debet in su%
